% system-model.tex -- Batch Aggregation Paper (RU-V4)

\section{System Model}

\subsection{Network Architecture}

The Basis Network Enterprise ZK Validium operates as a permissioned
application-specific node serving a single enterprise tenant. Transactions
originate from enterprise systems (PLASMA industrial maintenance or Trace ERP)
and represent state transitions in a Sparse Merkle Tree (SMT) of enterprise
state~\cite{baylina2019smt}. The node collects transactions, aggregates them
into fixed-size batches, generates Groth16 ZK proofs via
Circom/SnarkJS~\cite{groth2016size, snarkjs}, and submits proof-state root
pairs to the Basis Network L1 on Avalanche for settlement.

The batch aggregation layer sits between transaction ingestion and proof
generation:
\begin{center}
\small
Enterprise Tx $\rightarrow$ [Queue + WAL] $\rightarrow$ [BatchAggregator] $\rightarrow$ [ZK Prover] $\rightarrow$ L1
\end{center}

\subsection{Trust Model}

The enterprise node is operated by a single entity with administrative
access. The trust model assumes:

\begin{itemize}
\item \textbf{Honest operator}: The enterprise operator does not intentionally
  corrupt the WAL or manipulate batch formation. Malicious operators are outside
  the threat model (they control the data source).
\item \textbf{Byzantine faults excluded}: The node is a single process on a
  single machine. There is no consensus protocol between nodes.
\item \textbf{Crash-stop failures}: The node may crash at any point due to
  hardware failure, power loss, out-of-memory conditions, or process termination.
  After crash, the node restarts and recovers from durable state.
\item \textbf{Durable storage reliability}: The WAL on disk survives crashes.
  Storage corruption (bit rot, sector failures) is excluded from the model.
\end{itemize}

\subsection{Safety Requirements}

The batch aggregation protocol must satisfy four invariants:

\begin{enumerate}
\item \textbf{NoLoss}: Every transaction that enters the system is eventually
  included in a processed batch. No transaction is silently dropped, even under
  crash-recovery sequences.
\item \textbf{NoDuplication}: No transaction appears in more than one batch or
  is processed more than once.
\item \textbf{FIFOOrdering}: Transactions are batched and processed in the
  exact order they were received. The concatenation of all processed and pending
  batches equals the WAL in arrival order.
\item \textbf{BatchSizeBound}: Every batch contains at most
  \texttt{BatchSizeThreshold} transactions, ensuring compatibility with the
  ZK circuit's constraint capacity.
\end{enumerate}

These requirements are formalized as TLA+ invariants in Section~\ref{sec:formal}.

\subsection{Performance Requirements}

Based on the enterprise deployment context (industrial traceability at
Ingenio Sancarlos, Colombia), the system must achieve:

\begin{itemize}
\item Throughput $\geq$ 100 tx/min sustained
\item Batch formation latency $<$ 5,000 ms
\item Zero transaction loss under crash recovery
\item Deterministic batch formation (same inputs produce identical outputs)
\end{itemize}
